\chapter{JavaScript: le basi}

\section{Carta d'identit\`a del linguaggio}

\begin{definizione}[Caratteristiche di JavaScript]
\begin{itemize}
  \item \`E un linguaggio \textbf{interpretato}: necessita di un motore di esecuzione
        (interprete).
  \item \`E \textbf{imperativo e strutturato}: programmi fatti di istruzioni, organizzabili
        in blocchi, con condizionali e cicli.
  \item \`E \textbf{debolmente tipato}: i dati vengono convertiti automaticamente da un tipo
        all'altro se il calcolo di un'espressione lo richiede (35 pu\`o diventare
        \texttt{"35"}, \texttt{false} pu\`o diventare 0).
  \item \`E \textbf{dinamicamente tipato}: il tipo di un'espressione \`e stabilito solo a
        run-time; non si dichiara in anticipo il tipo di variabili, parametri o valori di
        ritorno.
  \item \`E \textbf{object-oriented basato su prototipi}: gli oggetti non hanno bisogno di
        classi per essere creati (un oggetto \`e una collezione di coppie propriet\`a/valore;
        un oggetto pu\`o fare da \emph{prototipo} a un altro, che ne eredita le propriet\`a;
        le classi esistono comunque, per semplificare la creazione di oggetti con le stesse
        propriet\`a e lo stesso prototipo).
  \item Supporta la \textbf{programmazione funzionale}: le funzioni sono \emph{tipi di dati}
        -- possono essere create dinamicamente, memorizzate in variabili, passate come
        parametri, assegnate a propriet\`a di oggetti.
\end{itemize}
\end{definizione}

\subsection{JavaScript o ECMAScript?}
JavaScript nasce come linguaggio di scripting per il web; dopo un periodo ``selvaggio'' in
cui ogni browser dava la propria interpretazione, si \`e imposta la standardizzazione
\textbf{ECMAScript} (ES): la specifica ufficiale, in costante evoluzione. ``JavaScript'' \`e
l'implementazione della specifica realizzata da Mozilla (Firefox); ormai tutti i browser
implementano correttamente ECMAScript. Le feature aggiunte nelle edizioni successive servono a
una programmazione pi\`u concisa e manutenibile, raramente introducono costrutti non
replicabili nelle versioni precedenti: per questo un programma per una versione recente pu\`o
essere \emph{tradotto automaticamente} in una versione precedente. Questa traduzione, detta
\textbf{transpilazione}, \`e diffusa per far girare i programmi anche su browser meno
aggiornati. L'ultima edizione che ha modificato sostanzialmente il linguaggio \`e stata
\textbf{ES2015 (ES6)}; il corso non usa nulla di posteriore a \textbf{ES2020 (ES11)}.

\subsection{Linguaggi interpretati, compilati, ibridi}
Le istruzioni sono eseguite da un \emph{agente di calcolo}, che decide il set di istruzioni
valide e il loro significato. L'agente di base \`e il processore (linguaggio macchina,
inutilizzabile direttamente). I linguaggi di alto livello si dividono in:
\begin{itemize}
  \item \textbf{compilati} (es. C): il programma viene tradotto \emph{una volta per tutte} in
        linguaggio macchina da un \emph{compilatore}; il risultato \`e un eseguibile,
        specifico per la famiglia di processori;
  \item \textbf{interpretati} (es. JavaScript): usano un agente di calcolo di livello pi\`u
        alto -- l'\emph{interprete} o \emph{motore di esecuzione} (lui s\`i compilato) -- che
        esegue direttamente le istruzioni. L'agente realizza una \textbf{macchina virtuale}:
        non esiste una macchina fisica che esegue quelle istruzioni, ma per chi programma \`e
        come se ci fosse;
  \item \textbf{ibridi} (es. Java): si compila in un linguaggio intermedio (Java ByteCode)
        che viene poi interpretato da un motore (la JVM). Lo schema pu\`o avere anche pi\`u
        livelli: si potrebbe inventare un linguaggio ``VoiceScript'' con un motore scritto in
        JavaScript, il quale a sua volta gira su un motore JavaScript, che invoca istruzioni
        macchina.
\end{itemize}

\section{Eseguire JavaScript nel browser}

Il \textbf{motore di esecuzione} realizza un \textbf{ambiente di esecuzione}: ``programma in
esecuzione'' + ``memoria di esecuzione''. \`E un ambiente \emph{vivo} (oggetti e variabili
aggiunti, tolti, modificati) e pu\`o essere corredato di \textbf{API} che offrono costrutti
per interagire con quell'ambiente. Nel browser: il motore crea una \emph{sandbox} separata
per ogni pagina, resettata al ricaricamento. L'ambiente contiene l'oggetto \textbf{document}
(il DOM della pagina), la \textbf{DOM API} per consultarlo e modificarlo, e la
\textbf{fetch API} per le richieste HTTP.

\subsection{Collegare uno script a una pagina}
\begin{lstlisting}[language=HTMLCSS]
<script type="application/javascript" src="path/to/myscript.js"></script>
\end{lstlisting}
Tipicamente in HEAD; lo script viene per\`o eseguito \emph{immediatamente}, senza aspettare il
caricamento della pagina -- un problema se interagisce con \texttt{document}, che potrebbe
essere ancora vuoto. Tre modi per rimandare l'esecuzione:
\begin{enumerate}
  \item spostare SCRIPT \emph{al fondo}, dopo la chiusura di BODY (ma prima di quella di HTML);
  \item usare l'attributo \textbf{defer};
  \item dichiarare \texttt{type="module"}: indica che il file \`e il \emph{punto d'ingresso}
        di un'applicazione su pi\`u moduli. \`E la soluzione pi\`u usata.
\end{enumerate}

\section{Tipi, variabili e strutture di controllo}

\subsection{I tipi base}
\texttt{number}, \texttt{boolean}, \texttt{string}, \texttt{object}, \texttt{undefined},
\texttt{null}, \texttt{bigint}, \texttt{symbol} (array e funzioni sono oggetti).
\begin{itemize}
  \item Nessuna distinzione interi/decimali: tutti \texttt{number}, floating point a 64 bit
        (come i \texttt{double} di Java); \texttt{bigint} per interi molto grandi.
  \item Le stringhe (codifica UTF-16) si definiscono con virgolette \texttt{"}, apici
        \texttt{'} (equivalenti, salvo poter usare l'uno dentro l'altro senza escape) o
        \textbf{backtick} \texttt{`}: queste ultime sono \emph{template literal} -- tutto ci\`o
        che sta fra i backtick fa parte della stringa, spazi e tabulazioni inclusi (e, come
        vedremo negli esempi, permettono l'interpolazione \texttt{\$\{espr\}}).
  \item \texttt{undefined} e \texttt{null}: ciascuno \`e un tipo con un unico valore.
        Entrambi rappresentano operazioni non andate a buon fine; \texttt{undefined} =
        ``nessun valore'' (variabile non assegnata, funzione che non restituisce nulla),
        \texttt{null} = ``oggetto nullo''. In pratica usati in modo intercambiabile: meglio
        sceglierne uno e attenersi, o avere ben chiaro quando si usa l'uno e quando l'altro.
\end{itemize}

\subsection{Variabili}
\begin{lstlisting}[language=JavaScript]
let anno = 2024;      // dichiarazione, inizializzazione opzionale
const pi = 3.14;      // va inizializzata subito, non riassegnabile
\end{lstlisting}
Le variabili \emph{non hanno tipo}: possono contenere valori di tipi diversi in momenti
diversi.
\begin{attenzione}[const: costante $\neq$ immutabile]
Il valore di una \texttt{const} \`e \emph{costante} ma non \emph{immutabile}: se contiene un
tipo-riferimento (oggetto, array), il riferimento non cambia (stessa identit\`a in memoria) ma
si possono aggiungere/togliere elementi all'array o modificare le propriet\`a dell'oggetto.
\end{attenzione}

\subsection{Operatori e uguaglianza}
Operatori aritmetici e booleani standard; \texttt{+} concatena anche stringhe (come Java).
Poich\'e JS \`e debolmente tipato, \texttt{12 == "12"} risulta \emph{vero}: per questo
esistono \texttt{===} e \texttt{!==}:
\begin{center}
\texttt{(a === b)} $\Leftrightarrow$ \texttt{(a == b \&\& typeof a == typeof b)}
\end{center}
\begin{attenzione}
\texttt{==} e \texttt{!=} sono fonte di innumerevoli bug (es. \texttt{undefined} convertito a
\texttt{false} o 0 maschera errori di calcolo): buona prassi usare \emph{solo}
\texttt{===} e \texttt{!==}.
\end{attenzione}

\subsection{Controllo di flusso}
\texttt{if}, \texttt{while}, \texttt{do\dots while} identici a C/Java. Tre varianti di
\texttt{for}:
\begin{itemize}
  \item \texttt{for(init; condition; increment)}: il for classico;
  \item \texttt{for (const elem \textbf{of} iterabile)}: itera sugli \emph{elementi} di una
        struttura iterabile (array, string sui caratteri, \texttt{Set}, \texttt{Map});
  \item \texttt{for (const key \textbf{in} oggetto)}: itera sulle \emph{propriet\`a} di un
        oggetto (\texttt{key} \`e il nome della propriet\`a come stringa).
\end{itemize}

\section{Array}

\begin{lstlisting}[language=JavaScript]
let amici = ["Antonio", "Paola", "Giorgio", "Andrea", "Ellen"];
let vuoto = [];
amici[2]  // lettura/scrittura per indice
\end{lstlisting}
\begin{itemize}
  \item Nessuna dimensione prestabilita: la lunghezza \`e ultima posizione occupata + 1.
  \item Possono essere \emph{sparsi}: assegnando la posizione 100 le posizioni 5--99 restano
        \texttt{undefined}.
  \item Possono mescolare tipi diversi.
\end{itemize}
Gli array sono \emph{oggetti} della classe \texttt{Array}, con propriet\`a e metodi:
\begin{itemize}
  \item \texttt{a.length}: la lunghezza;
  \item \texttt{a.pop()}: restituisce, \emph{rimuovendolo}, l'ultimo elemento;
  \item \texttt{a.push(v)}: aggiunge \texttt{v} in coda;
  \item \texttt{a.join(sep)}: stringa con gli elementi concatenati, separati da \texttt{sep};
  \item \texttt{a.slice(start, end)}: \emph{nuovo} array con gli elementi da \texttt{start}
        inclusa a \texttt{end} esclusa;
  \item \texttt{a.splice(start, delNum, add1, \dots, addN)}: \emph{modifica} l'array attuale
        da \texttt{start}: cancella \texttt{delNum} elementi e inserisce i nuovi; con solo
        cancellazione (nessun elemento elencato) o solo aggiunta (\texttt{delNum} = 0).
\end{itemize}

\section{Oggetti base}

\begin{definizione}[Object literal]
Un oggetto \`e una collezione di coppie propriet\`a/valore. La forma pi\`u semplice \`e
l'\textbf{object literal}: propriet\`a \texttt{prop: valore} separate da virgole, fra graffe.
\begin{lstlisting}[language=JavaScript]
const libro = {
  id: 8,
  titolo: "Pedagogia degli oppressi",
  autore: "Paulo Freire",
  stato: "borrowed"
}
\end{lstlisting}
I nomi devono essere ben formati (come le variabili: alfanumerici, \texttt{\$},
\texttt{\_}; il primo carattere non pu\`o essere una cifra).
\end{definizione}

Accesso in lettura/scrittura con la \textbf{dot notation} (\texttt{libro.titolo}) o la
\textbf{square bracket notation} (\texttt{libro["titolo"]}). Si possono assegnare nuove
propriet\`a dopo la creazione (\texttt{libro.editore = "Neri Pozza";}); con le quadre si
possono usare nomi non validi per la dot notation
(\texttt{libro["ISBN-13"] = "\dots";} -- contiene un trattino). Si possono anche cancellare:
\texttt{delete libro.stato;}.

\section{Classi (introduzione)}

\begin{lstlisting}[language=JavaScript]
class Studente {
  constructor(nome) {
    this.nome = nome;
    this.votiEsami = [];
  }
  aggiungiVoto(voto) { this.votiEsami.push(voto); }
}
const s = new Studente("Saverio");
s.aggiungiVoto(28);
\end{lstlisting}
\begin{itemize}
  \item Il \textbf{costruttore} (\texttt{constructor}) inizializza le propriet\`a; l'oggetto
        in corso di inizializzazione \`e \texttt{this}.
  \item I \textbf{metodi} sono funzioni possedute da tutti gli oggetti della classe;
        \texttt{this} si riferisce all'oggetto su cui il metodo \`e invocato.
  \item Non c'\`e una sezione che elenca le propriet\`a: gli oggetti possono
        acquisirne/perderne dopo la creazione, e senza tipi l'elenco non darebbe vantaggi.
  \item Propriet\`a e metodi sono \emph{tutti pubblici}; per l'``uso interno'' si usa la
        convenzione dei nomi con underscore iniziale (non \`e un vincolo reale).
  \item \textbf{Ereditariet\`a}: \texttt{class Laureato extends Studente \{\dots\}};
        il costruttore della sottoclasse deve invocare \emph{per prima cosa} quello della
        superclasse con \texttt{super(\dots)}, poi imposta le propriet\`a aggiuntive. Gli
        oggetti della sottoclasse ereditano propriet\`a e metodi della superclasse.
\end{itemize}

\section{Funzioni}

Tre modi di dichiararle:
\begin{lstlisting}[language=JavaScript]
function myfun(x, y, z) { ...; return valore; }        // named
const noname = function (x, y, z) { ...; return v; };  // anonymous
const arfun = (x, y, z) => { ...; return v; };         // arrow
const oneline = (x, y, z) => valore;                   // arrow, corpo-espressione
\end{lstlisting}
\begin{attenzione}
Invocando una funzione \emph{non} \`e obbligatorio passare tutti i parametri: quelli mancanti
valgono \texttt{undefined}. Una funzione senza \texttt{return} esplicita restituisce
\texttt{undefined}. Quindi un'invocazione sbagliata o un return dimenticato producono errori
che ``saltano fuori'' in punti molto diversi del programma.
\end{attenzione}

\textbf{Funzioni come dati.} Le funzioni si possono dichiarare ovunque (visibili nello scope
di creazione); si possono assegnare a variabili e a propriet\`a di oggetti. \`E importante
distinguere:
\begin{itemize}
  \item la \emph{funzione in quanto oggetto} (indicata col suo nome, o col nome della
        variabile che la contiene): la funzione \`e un dato;
  \item l'\emph{invocazione} (nome + parentesi tonde + parametri): un'espressione con un
        risultato (che pu\`o essere \texttt{undefined}\dots\ o un'altra funzione!).
\end{itemize}
Una funzione pu\`o essere passata come parametro a un'altra e invocata al suo interno col
nome del parametro formale:
\begin{lstlisting}[language=JavaScript]
function applica(testo, analisi) { return analisi(testo); }
let contaparole = s => s.split(' ').length;
let n = applica("c'era una volta", contaparole);
\end{lstlisting}

\begin{attenzione}[Le propriet\`a-funzione non sono metodi]
Assegnando una funzione a una propriet\`a di un oggetto \emph{non} si ottiene automaticamente
un metodo: se la funzione \`e una \emph{arrow function}, durante l'invocazione \textbf{non \`e
attivo il binding di \texttt{this}} (this non viene ``vincolato'' all'oggetto proprietario).
Il tema del binding \`e approfondito nel capitolo sui prototipi.
\end{attenzione}

\section{Scope e closure}

\begin{definizione}[Scope]
Ogni blocco di istruzioni fra graffe istituisce uno \textbf{scope} (ambito di visibilit\`a):
blocchi \texttt{if}/\texttt{else}, \texttt{while}/\texttt{for}, corpi di funzione, o anche
blocchi arbitrari. Gli scope possono essere \emph{innestati}; il pi\`u esterno (l'intero
script) \`e lo \textbf{scope globale}. La visibilit\`a di variabili e nomi di funzione \`e
limitata allo scope di definizione e ai suoi discendenti.
\end{definizione}

\begin{definizione}[Closure]
Le variabili di uno scope \textbf{sopravvivono alla sua chiusura} fino a quando esiste uno
scope interno attivo che le utilizza: questo meccanismo \`e detto \textbf{closure}.
\`E tipico dei linguaggi funzionali ed \`e necessario perch\'e le funzioni possano essere
``elementi di prima classe''. In altri termini: la graffa chiusa segna la conclusione
\emph{statica} dello scope, ma l'esecuzione dell'ultima istruzione non ne segna la fine della
\emph{vita dinamica}.
\end{definizione}

\begin{esempio}
\begin{lstlisting}[language=JavaScript]
function functionCreator() {
  let somearray = ["1..", "2..", "3..", "Prova"];
  let somefun = (testo) => {
    for (let elem of somearray)
      console.log(`${elem} --> ${testo}`);
  }
  let update = (s) => { somearray.push(s); console.log("Aggiornato!"); }
  update("Via!");
  return somefun;
}
const risultato = functionCreator();
risultato("Eccomi!");  // somearray e' ancora "vivo"!
\end{lstlisting}
\texttt{somearray} \`e creato nello scope di \texttt{functionCreator} e usato da una funzione
interna, che viene restituita e invocata \emph{dopo} che \texttt{functionCreator} \`e
terminata: grazie alla closure, \texttt{somearray} sopravvive al suo creatore.
\end{esempio}

\section{Array e programmazione funzionale}

Il fatto che le funzioni siano dati passabili come parametri permette uno stile
\textbf{funzionale} e dunque \emph{dichiarativo}: non si impartiscono istruzioni
algoritmiche su \emph{come} ottenere il risultato, si dichiara \emph{cosa} si vuole ottenere.
\begin{esempio}[map: imperativo vs dichiarativo]
\begin{lstlisting}[language=JavaScript]
// imperativo
const numeri = [1,5,8,12,13];
const quadrati = [];
for (let n of numeri) quadrati.push(n*n);
// dichiarativo
const quadrati2 = numeri.map(x => x * x);
\end{lstlisting}
\texttt{map} prende una funzione di mappatura (trasforma un singolo elemento) e restituisce
un \emph{nuovo} array con tutti gli elementi mappati.
\end{esempio}

Altri metodi ``dichiarativi'' degli array:
\begin{itemize}
  \item \texttt{a.forEach(fun)}: esegue \texttt{fun} per ogni elemento (il valore di ritorno
        di \texttt{fun} \`e ignorato);
  \item \texttt{a.find(searchFun)}: primo elemento per cui \texttt{searchFun} (che restituisce
        boolean) \`e vera; esiste anche \texttt{findLast};
  \item \texttt{a.findIndex(searchFun)}: come find ma restituisce la \emph{posizione}
        ($-1$ se nessun elemento soddisfa); anche \texttt{findLastIndex};
  \item \texttt{a.every(testFun)}: true se \texttt{testFun} \`e vera per \emph{tutti};
  \item \texttt{a.some(testFun)}: true se vera per \emph{almeno uno};
  \item \texttt{a.reduce(accumulatore, valoreIniziale)}: ``accumula'' i valori in un unico
        risultato; l'accumulatore \`e una funzione (valore accumulato, prossimo valore)
        $\rightarrow$ risultato dello step. Esempio, la somma:
        \texttt{numeri.reduce((somma, v) => somma + v, 0);}
\end{itemize}

\section{Moduli: organizzare il codice su pi\`u file}

In JavaScript (a differenza di Java) il codice pu\`o essere diviso su pi\`u file in modo
arbitrario. In teoria: uno script non viene compilato n\'e controllato prima dell'esecuzione,
quindi pu\`o usare funzioni/oggetti definiti altrove \emph{purch\'e} quei file siano stati
eseguiti prima e l'ambiente sia popolato. Il modo rudimentale (ormai quasi solo in
applicazioni molto vecchie) \`e includere gli script nell'HTML \emph{nell'ordine giusto}: non
regge per applicazioni complesse che usano estensivamente librerie di terze parti, dove
\`e cruciale che le \textbf{dipendenze} siano chiare.

\begin{definizione}[Moduli]
Ogni file JavaScript \`e un \textbf{modulo}. Un modulo pu\`o \textbf{importare} ed
\textbf{esportare} oggetti, variabili, funzioni e classi.
\end{definizione}

\subsection{Export}
\begin{itemize}
  \item \textbf{standard export}: si premette \texttt{export} alla dichiarazione:
\begin{lstlisting}[language=JavaScript]
export function pincopallino() {...}
export let nomeVar = ...
export class UnaClasse {...}
\end{lstlisting}
  \item \textbf{default export}: quando un modulo esporta \emph{una sola} cosa; si premette
        \texttt{export default}. Ci\`o che non viene esportato resta \emph{privato} del
        modulo -- \`e anche un modo di proteggerne i meccanismi interni.
\end{itemize}

\subsection{Import}
\begin{itemize}
  \item \textbf{default import}: \texttt{import nomeQualsiasi from "altroscript.js"} --
        il nome \`e a scelta, tanto l'elemento esportato \`e uno solo;
  \item \textbf{standard import di base}:
        \texttt{import \{cosa1, cosa2\} from "altroscript.js"} -- si usano i nomi originali;
  \item \textbf{con renaming}:
        \texttt{import \{cosa1 as thing1\} from "altroscript.js"} -- si possono rinominare
        solo alcuni elementi;
  \item \textbf{con namespace}: \texttt{import * as Qualcosa from "altroscript.js"} --
        importa tutto, con prefisso: \texttt{Qualcosa.x}.
\end{itemize}
